\documentclass[UTF8,11pt,a4paper]{ctexart}
\newcommand{\DocTitle}{游戏与规则}
\newcommand{\DocSubtitle}{从开局采购到每回合结算}
% Shared layout for the CW2 companion documents.
\usepackage[margin=2.1cm,headheight=16pt]{geometry}
\usepackage{amsmath,amssymb,booktabs,array,tabularx,ragged2e}
\usepackage{enumitem,xcolor,fancyhdr,listings,tcolorbox}
\tcbuselibrary{breakable}
\usepackage{hyperref}
\usepackage{bookmark}

\newcommand{\DocVersion}{3.0}
\newcommand{\EngineVersion}{1.1}
\newcommand{\DocDate}{2026-11-23}
\definecolor{ink}{HTML}{1F2937}
\definecolor{accent}{HTML}{92400E}
\definecolor{muted}{HTML}{4B5563}
\definecolor{linkink}{HTML}{174678}
\definecolor{codebg}{HTML}{F5F7FA}
\hypersetup{unicode,colorlinks=true,linkcolor=linkink,urlcolor=linkink,
  filecolor=linkink,pdftitle={SOE104 CW2 \DocTitle\space v\DocVersion},
  pdfauthor={Group A--E}}
\ctexset{section={format={\Large\bfseries\color{accent}}},
  subsection={format={\large\bfseries\color{ink}}}}
\pagestyle{fancy}
\fancyhf{}
\fancyhead[L]{\small\color{muted}SOE104 CW2 · Oregon Trail}
\fancyhead[R]{\small\color{muted}\DocTitle\ v\DocVersion}
\fancyfoot[C]{\small\color{muted}\thepage}
\renewcommand{\headrulewidth}{0.4pt}
\setlength{\parskip}{0.35em}
\setlength{\emergencystretch}{1.5em}
\renewcommand{\arraystretch}{1.23}
\setlist{itemsep=3pt,topsep=3pt,parsep=0pt,leftmargin=1.7em}
\newcolumntype{Y}{>{\RaggedRight\arraybackslash}X}
\newcolumntype{L}[1]{>{\RaggedRight\arraybackslash}p{#1}}
\newtcolorbox{keybox}[1]{breakable,title={#1},colback=orange!4,
  colframe=accent,boxrule=0.55pt,arc=2pt,fonttitle=\bfseries}
\newtcolorbox{warnbox}[1]{breakable,title={#1},colback=red!3,
  colframe=red!45!black,boxrule=0.55pt,arc=2pt,fonttitle=\bfseries}
\lstset{basicstyle=\small\ttfamily,backgroundcolor=\color{codebg},
  xleftmargin=0.9em,frame=l,rulecolor=\color{linkink!45},
  breaklines=true,columns=fullflexible,keepspaces=true,showstringspaces=false,
  aboveskip=0.65em,belowskip=0.65em}
\newcommand{\code}[1]{\texttt{#1}}
\newcommand{\source}[1]{\textup{\small\textcolor{muted}{[#1]}}}
\newcommand{\docref}[2]{\href{./#1.pdf}{#2}}
\renewcommand{\maketitle}{%
  \begin{center}
    {\LARGE\bfseries CW2 \DocTitle\par}
    \vspace{0.35em}{\large\DocSubtitle\par}
    \vspace{0.75em}{\small\color{muted}Group A--E · 整理：A
      \quad\DocDate\quad 文档 v\DocVersion\par}
    \vspace{0.2em}{\small\color{muted}Engine / Marking v\EngineVersion\par}
  \end{center}
  \vspace{0.4em}}


\begin{document}
\maketitle
\thispagestyle{fancy}

\begin{keybox}{你控制哪四件事}
团队提交一组 Java 类。课程引擎负责随机事件、资源结算、回放和评分。
我们决定开局买什么、每回合怎么走和吃、怎么应对事件，以及成员各自做什么。
\end{keybox}

\section{旅程与回合}

目标是让 3--5 名游戏成员在 120 回合内到达 2000 英里。旅程在到达、用完回合或全队死亡时结束。
牛降到 0 会失去配速行进，回合与喂食仍继续。
平台使用同样的三个种子评估各队，比较两个版本时应按种子配对。

\begin{center}\small
\begin{tabularx}{\linewidth}{@{}L{2.5cm}L{4.7cm}Y@{}}
\toprule
时点 & 你的入口 & 当时能读到什么 \\
\midrule
开局采购 & \code{buySupplies()} & 返回购买清单，引擎按顺序结算。 \\
回合开始 & \code{decide(PartyState)} & 开头的里程、成员、库存，返回配速和口粮。 \\
行进、喂食后 & \code{resolve(Event)} & 当前事件和严重度，市场事件有报价。 \\
事件结算后 & \code{act(Situation)} & 成员行动时的资源和事件信息，存活成员按队列顺序行动。 \\
\bottomrule
\end{tabularx}
\end{center}
随后是回合结算和生存检查。事件扣粮发生在本轮喂食之后，会挤压下一餐与当前行动的粮食预算。
\source{ER} \source{API}

\subsection{配速与运力}
配速里程为基础里程乘 $\min(1,O/4)$，$O$ 为牛数。
\code{GRUELING} 每人损失 4 HP，并有 15\% 概率损失一头牛。
有活 Driver 时分别降为 2 HP 和 7\%。 \source{ER}
\begin{center}\small
\begin{tabularx}{\linewidth}{@{}L{3.0cm}YYY@{}}
\toprule
牛数 & SLOW & NORMAL & GRUELING \\
\midrule
4 或更多 & 12 & 20 & 28 \\
3 & 9 & 15 & 21 \\
2 & 6 & 10 & 14 \\
1 & 3 & 5 & 7 \\
0 & 0 & 0 & 0 \\
\bottomrule
\end{tabularx}
\end{center}
四头牛用 \code{NORMAL}、只算配速需 100 回合。
三头牛用 \code{NORMAL}、只算配速需 134 回合，但 Scout 每回合可另贡献 6 英里，
还有提速和途中补牛方案。是否可达要看剩余里程、动作分工和绕道损失。
\source{DER}

\clearpage
\section{供餐、动作与职业}

成员先吃，牛随后每头吃 2 食物。口粮为 \code{BARE} 1、\code{MEAGER} 3、\code{FULL} 6 食物/人。
充分供餐时三档分别是每人 $-1$、$0$、最多 $+1$ HP。
四人四牛用 \code{MEAGER} 时，喂饱成员需 12 食物，喂饱全车需 20。
成员吃不饱才计饥饿，连续饥饿每人依次扣 5、8、12、17、23 HP，此后保持 23。
一次吃饱重置计数。健康上限 100，到 0 即死亡。 \source{TU} \source{ER} \source{DER}

\begin{center}\small
\begin{tabularx}{\linewidth}{@{}L{2.5cm}Y Y@{}}
\toprule
动作 & 基础效果 & 职业或使用限制 \\
\midrule
\code{HUNT} & 2 弹药换最多 7 食物 & Hunter 最多 10。全车食物上限 500。 \\
\code{HEAL} & 3 食物，自身最多恢复 3 HP & Doctor 治疗最伤的活成员，最多 5 HP。挨饿当轮无效。 \\
\code{REST} & 自身最多恢复 2 HP & 挨饿当轮无效。 \\
\code{SCOUT} & 本轮增加 3 英里 & Scout 增加 6 英里。 \\
\code{MOVE} & 本轮增加 5 英里，自身损失 1 HP & Driver 也掉 1 HP。 \\
\code{DEFEND} & 受威胁时每名活成员最多恢复 3 HP & 整队每轮至多触发一次。 \\
\bottomrule
\end{tabularx}
\end{center}
活着的 Doctor 将全队受到的\textbf{事件伤害}减半。活着的 Mechanic 可免费处理
\code{BREAKDOWN} 的 \code{SAFE} 修理。Driver 的光环作用于艰难配速。
这三类光环各自不叠加。 \source{TU} \source{ER}

\subsection{读回放时，按这个顺序核对}

例如四人四牛在回合开始时有 30 食物，普通口粮（\code{MEAGER}）喂完车队后剩 10。
随后遇到严重度 4 的失窃，选择 \code{SAFE} 扣 9 食物，车里只剩 1。
如果猎人有 2 发弹药，本轮狩猎后可把粮补到 11。这个回合队员已经吃饱，
但下次要继续用普通口粮，四人至少要有 12 食物，连牛一起则要 20。
回放上看到的「仓里只剩一点粮」，要和\textbf{发生在哪个阶段}一起读。 \source{DER}

\begin{keybox}{三个容易看错的动作结果}
\begin{itemize}
  \item 499 食物时狩猎，食品仓位只够再加 1，弹药仍扣 2 发。
  \item 成员吃不饱的回合，\code{HEAL} 和 \code{REST} 无效，狩猎仍能补给。
  \item 受威胁且需要防守时，一人使用 \code{DEFEND} 即可，队员各自最多恢复到 100 HP。
\end{itemize}
\end{keybox}

\clearpage
\section{事件与交易}

1.1 公告给出的事件率从旅途前期约 20\% 上升到后期约 56\%，
完整的新版概率曲线尚待核验。表中 $s$ 是事件传来的严重度（1--10）。
伤害为减伤前每人值。仓位限制获得的食物，现有库存限制实际扣粮。
河流与袭击的付费选项在钱不够时会转为 \code{AVOID}。 \source{TU} \source{ER} \source{API}

\begin{center}\small
\begin{tabularx}{\linewidth}{@{}L{2.55cm}Y Y L{2.25cm}@{}}
\toprule
事件 & SAFE & RISKY & AVOID \\
\midrule
\code{RIVER} & 付 $20+s$ 美元 &
$s/10$ 概率掉 10 HP 和 20 粮，另有 20\% 失牛检定 & 损失 5 英里 \\
\code{STORM} & 掉 $s$ HP & 相同伤害，另有 15\% 失牛风险 & 损失 10 英里 \\
\code{BREAKDOWN} & 有机械师免零件，否则用 1 件；无零件时按绕行处理 &
$s/10$ 概率掉 5 HP 并失 1 零件 & 损失 15 英里 \\
\code{ILLNESS} & 随严重度变化，精确成本待确认 &
$s/10$ 概率掉 $s+5$ HP & 掉 $s$ HP \\
\code{ATTACK} & 付 $20+s$ 美元，失 10 粮 & 掉 $s$ HP & 损失 15 英里 \\
\code{HUNT\_OPP} & 用 2 发，最多得 20 粮 & 用 4 发，最多得 50 粮；$s/10$ 概率掉 5 HP & 放弃 \\
\code{BAD\_TRAIL} & 损失 $6+s$ 英里；旧版扣 5 粮是否保留待核 &
$s/10$ 概率失败并掉 $s$ HP，失败时另有 15\% 失牛风险；否则加 5 英里 & 损失 20 英里 \\
\code{THEFT} & 失 $5+s$ 粮 & 各半概率掉 5 HP，或失 20 粮和 \$20 & 失 \$15 \\
\bottomrule
\end{tabularx}
\end{center}

\noindent 河流和坏路的失牛检定细节仍需通过回放核验。

\subsection{商店与交易站}
开局商店只买不卖，钱包 \$2,000。每单位粮、弹药、零件、牛分别为
\$2、\$3、\$30、\$80。经典篮子为 500 粮、80 弹、2 零件、5 牛，余 \$300。
采购清单按顺序扣钱，超出余额或仓位时会裁剪。 \source{ER}

途中遇到 \code{TRADER} 时可买、卖、议价或离开，每次只做一笔交易。
严重度 1--3 为便宜档、4--6 为正常档、7--10 为贵档。按事件给的报价算账。
买入最多 50 单位，受钱包和仓位限制；卖出不能超过库存。议价失败则不成交。
例如仓里有 490 粮、现金 \$12，挂牌粮价 \$2，申请买 50 也只买到 6。 \source{ER} \source{DER}

\medskip
\noindent 事件成本与可达性计算见\docref{04_decision_model}{《04 决策模型》}，
分数和提交见\docref{02_grading}{《02 评分与提交》}。
\source{TU} 指 Trail Update 1.1，\source{ER} 指未被公告替换的 Engine Rules Spec，
\source{API} 指 API Reference，\source{DER} 指由规则推导的例子。

\end{document}
